% SPDX-License-Identifier: CC-BY-4.0
\documentclass[11pt,letterpaper]{article}

\usepackage[T1]{fontenc}
\usepackage{amsmath,amssymb,amsthm}
\usepackage{mathtools}
\usepackage{microtype}
\usepackage{geometry}
\usepackage{hyperref}
\usepackage{enumitem}
\usepackage{booktabs}
\usepackage{array}

\geometry{margin=1in}
\hypersetup{colorlinks=true,linkcolor=blue,citecolor=blue,urlcolor=blue}

\newtheorem{theorem}{Theorem}[section]
\newtheorem{lemma}[theorem]{Lemma}
\newtheorem{proposition}[theorem]{Proposition}
\newtheorem{corollary}[theorem]{Corollary}
\theoremstyle{definition}
\newtheorem{definition}[theorem]{Definition}
\newtheorem{example}[theorem]{Example}
\theoremstyle{remark}
\newtheorem{remark}[theorem]{Remark}

\newcommand{\Z}{\mathbb{Z}}
\newcommand{\N}{\mathbb{N}}
\newcommand{\Q}{\mathbb{Q}}
\newcommand{\Map}{\mathcal{T}}
\newcommand{\Iter}{\Map^{\langle k\rangle}}

\title{Multiplicatively Independent Homogeneous Raw Maps}
\author{Ascedete et Ambulate}
\date{\today}

\begin{document}
\maketitle

\begin{abstract}
For integers $q,D\geq 2$, consider the homogeneous raw map that divides by
$D$ on divisible states and multiplies by $q$ otherwise.  If no positive
power of $q$ equals a positive power of $D$, every positive orbit is
injective and tends to infinity.  The proof retains the raw clock and every
branch guard.  A complete factor-free test for the power relation supplies
finite executable certificates without replacing the universal proof.
\end{abstract}

\section{Statement and conventions}

\begin{definition}[Positive homogeneous raw map]
Let $q,D\in\N$ satisfy $q,D\geq 2$.  For $n\in\N_{>0}$, define
\[
 \Map_{q,D}(n)=
 \begin{cases}
  n/D,&D\mid n,\\
  qn,&D\nmid n.
 \end{cases}
\]
The division branch has precedence.  Iteration is at the raw one-instruction
clock; no run of divisions is compressed.
\end{definition}

\begin{definition}[Multiplicative independence]
The pair $(q,D)$ is multiplicatively independent if
\[
 q^a\neq D^b
 \qquad\text{for every }a,b\in\N_{>0}.
\]
This is stronger than neither coprimality nor distinctness.  For example,
$6$ and $4$ are independent although they have a common prime factor.
\end{definition}

\begin{theorem}[Independent homogeneous escape]\label{thm:escape}
Let $q,D\geq 2$ be multiplicatively independent.  For every $n\geq 1$,
the raw orbit $k\mapsto\Iter_{q,D}(n)$ is injective and
\[
 \Iter_{q,D}(n)\longrightarrow\infty
 \qquad (k\longrightarrow\infty).
\]
Equivalently, for every $B\in\N$ there is $K$ such that
$\Iter_{q,D}(n)>B$ for all $k\geq K$.
\end{theorem}

The conclusion is genuine convergence to infinity, not merely unboundedness
and not merely the absence of cycles.

\section{Guarded word calculation}

Let a finite legal raw word $w$ contain $A(w)$ multiplication instructions
and $E(w)$ division instructions.  Thus
\[
 A(w)+E(w)=|w|.
\]
The order of these instructions controls legality, but in the homogeneous
case it does not change the relaxed monomial represented by the word.

\begin{lemma}[Exact homogeneous word law]\label{lem:word}
If a legal word $w$ carries a positive integer $x$ to $y$, then
\[
 D^{E(w)}y=q^{A(w)}x.
\]
In particular, every guard occurring in $w$ is retained in the hypothesis.
\end{lemma}

\begin{proof}
Associate to a multiplication instruction the rational affine data
$(P,B,Q)=(q,0,1)$ and to a division instruction the data $(1,0,D)$.
Chronological affine composition multiplies the numerators and denominators
and leaves the offset zero.  Induction on $w$ therefore gives
\[
 (P(w),B(w),Q(w))=(q^{A(w)},0,D^{E(w)}).
\]
Legal execution identifies $y$ with
$(P(w)x+B(w))/Q(w)$, which gives the displayed integer equation.  No
cancellation or converse legality assertion is used.
\end{proof}

\begin{lemma}[No nonempty positive return]\label{lem:noreturn}
Under the assumptions of Theorem~\ref{thm:escape}, no nonempty legal word
can carry a positive integer $x$ back to $x$.
\end{lemma}

\begin{proof}
Suppose a nonempty legal word returns $x>0$ to itself.  By
Lemma~\ref{lem:word},
\[
 (D^{E(w)}-q^{A(w)})x=0,
\]
so $D^{E(w)}=q^{A(w)}$.  Both instruction counts are positive.  Indeed, if
$A(w)=0$, then $D^{E(w)}=1$; because $D\geq2$, this forces $E(w)=0$.
Similarly, $E(w)=0$ forces $A(w)=0$.  Either conclusion contradicts
$|w|=A(w)+E(w)>0$.  Hence $A(w),E(w)>0$, and their equality of powers
contradicts multiplicative independence.
\end{proof}

\section{Recovery of full escape}

\begin{lemma}[Positivity]
Every raw iterate of a positive state is positive.
\end{lemma}

\begin{proof}
A multiplication step sends $x>0$ to $qx>0$.  If $D\mid x$, then
$x\geq D$ and the division step sends $x$ to the positive integer $x/D$.
Induction on the raw clock proves the claim.
\end{proof}

\begin{proposition}[Raw-time injectivity]\label{prop:injective}
Under multiplicative independence, the orbit of every positive seed is
injective.
\end{proposition}

\begin{proof}
Suppose the states at raw times $i<j$ coincide.  Determinism makes the
segment of length $j-i>0$ a legal return word at the positive intermediate
state.  This contradicts Lemma~\ref{lem:noreturn}.
\end{proof}

\begin{proof}[Proof of Theorem~\ref{thm:escape}]
Proposition~\ref{prop:injective} shows that the orbit is an injective
sequence in $\N_{>0}$.  Fix $B\in\N$.  At most $B$ orbit terms can lie in
the finite set $\{1,\ldots,B\}$.  Their time indices form a finite set, so
a time $K$ exceeds every one of them.  Consequently
$\Iter_{q,D}(n)>B$ for every $k\geq K$.  This is precisely convergence to
infinity.
\end{proof}

This last finite-set argument is the closing resource: a bounded interval
has only finitely many available values, and raw-time injectivity spends each
value at most once.  It upgrades the no-return obstruction to the required
full-sequence limit.

\section{A complete factor-free dependence test}

The executable checker compares the increasing lists
\[
 q,q^2,\ldots,q^{L_D}
 \quad\text{and}\quad
 D,D^2,\ldots,D^{L_q},
\]
where $L_D$ and $L_q$ may be taken as the binary bit lengths of $D$ and $q$.
At each comparison it advances the smaller power.  Thus it either returns
an exact equality $q^a=D^b$ or certifies that the two finite lists are
disjoint.

\begin{proposition}[Completeness of the finite test]\label{prop:test}
Positive powers of $q$ and $D$ coincide if and only if they coincide within
those two finite lists.
\end{proposition}

\begin{proof}
Only completeness needs proof.  Suppose $q^a=D^b$ for some positive
$a,b$, and divide $(a,b)$ by their greatest common divisor.  Unique prime
factorization, in its standard proportional-exponent form
\cite[Ch.~I]{HardyWright08}, gives an integer $t\geq2$ and coprime positive
integers $a_0,b_0$ such that
\[
 q=t^{b_0},\qquad D=t^{a_0}.
\]
Therefore $2^{a_0}\leq D$ and $2^{b_0}\leq q$.  Hence $a_0$ is no larger
than the binary bit length of $D$, and $b_0$ is no larger than that of $q$.
The finite merge must encounter the equality
$q^{a_0}=D^{b_0}$.  Soundness is immediate because every reported equality
is evaluated with exact integers.
\end{proof}

The test uses no factorization.  For example, it finds
\[
 (2^{100})^{101}=(2^{101})^{100}
\]
by exact power comparison.  Its output decides whether
Theorem~\ref{thm:escape} applies; it does not estimate an orbit.

\section{Example and boundaries}

\begin{example}
For $(q,D)=(6,4)$, the prime $3$ divides every $6^a$ and no $4^b$, so the
pair is independent.  Starting from $1$, the orbit begins
\[
 1,6,36,9,54,324,81,486,2916,729,\ldots.
\]
The decreases at division times show why monotonicity is unavailable.
Nevertheless, Theorem~\ref{thm:escape} proves that no state repeats and that
the complete sequence tends to infinity.
\end{example}

\begin{remark}[Excluded cases]
The theorem does not classify multiplicatively dependent pairs.  It also
does not cover $q=1$, $D=1$, zero or negative states, nonzero offsets, or
signed divisors.  These are separate dynamical regimes; none is inferred
from the independent positive theorem.
\end{remark}

\section{Formal and executable correspondence}

The Lean development proves the following chain:
\begin{enumerate}[label=(\roman*)]
 \item the two letter counts partition every raw word;
 \item the exact homogeneous word data are
       $(q^{A(w)},0,D^{E(w)})$;
 \item a nonempty positive guarded return contradicts independence;
 \item the natural positive map agrees at every raw time with the
       foundational signed-integer raw map on natural casts;
 \item every orbit is positive and injective;
 \item injectivity implies convergence to the top filter on $\N$.
\end{enumerate}

The arbitrary-precision executable checker independently verifies the
monomial invariant
\[
 x_jD^{E_j}=x_0q^{A_j}
\]
along bounded prefixes and checks that no positive prefix returns to its
seed.  Its exhaustive test covers $q,D\in\{2,\ldots,32\}$, seeds
$1\leq n\leq30$, and fifteen raw steps for every independent pair: exactly
$27060$ certificates and $432960$ endpoint-inclusive states.  A separate
finite test compares the power-relation detector with a redundant exact
search on all $q,D\in\{2,\ldots,128\}$.  These finite checks validate the
implementation; the universal conclusion is supplied by the Lean proof.

\bibliographystyle{amsalpha}
\bibliography{Homogeneous}

\end{document}
